\capitulo{2}{Objetivos del proyecto}

El presente Trabajo Fin de Grado tiene como objetivo principal el desarrollo de
una aplicación capaz de realizar análisis temáticos cualitativos sobre los
textos contenidos en las memorias de los \acrfull{tfg} en Ingeniería
Informática. Las memorias, \textit{issues}, \textit{commits}, documentos
descriptivos del \acrshort{tfg} y código fuente se encuentran disponibles en
repositorios de software almacenados en la plataforma GitHub.

Este proyecto busca aplicar técnicas de procesamiento del lenguaje natural y
modelado de temas para facilitar la identificación de tendencias, temas
recurrentes y áreas de investigación dentro del ámbito académico.

\section{Objetivo general}

Diseñar e implementar una aplicación que permita realizar análisis temáticos
automáticos sobre la documentación disponible en el repositorio del
\acrshort{tfg}, como \textit{readmes}, \textit{issues}, \textit{commits} y
memorias en formato \LaTeX{}, utilizando para esto diversos algoritmos de
modelado de temas y evaluando su rendimiento en términos de coherencia y
eficiencia.

\section{Objetivos específicos}

Para alcanzar el objetivo general, se plantean los siguientes objetivos
específicos:

\begin{itemize}
    \item \textbf{Recopilar y preparar los datos:} Crear un conjunto de datos utilizando la información de los  \acrshort{tfg} disponibles públicamente en GitHub, realizando el preprocesamiento necesario de los textos.

    \item \textbf{Implementar diferentes algoritmos de modelado de temas:} Investigar, aplicar y comparar métodos como \acrshort{lda}, \textbf{Top2Vec}, \textbf{BERTopic} y \textbf{FASTopic}.

    \item \textbf{Evaluar el rendimiento de los modelos:} Medir y comparar la calidad de los temas generados utilizando métricas de coherencia temática.

    \item \textbf{Desarrollar una interfaz visual:} Diseñar una aplicación que permita la visualización de los resultados obtenidos.

    \item \textbf{Fomentar la gestión del conocimiento académico:} Promover el acceso y la reutilización del conocimiento generado por estudiantes en cursos anteriores, apoyando la investigación y la docencia.

    \item \textbf{Diseño de arquitectura:} Diseñar e implementar una arquitectura de software modular que separe el motor de análisis de la capa de servicio, demostrando principios de Ingeniería de Software.

    \item \textbf{Gestión:} Utilizar herramientas y metodologías ágiles para planificar y implementar el proyecto. Utilizando Zube.io y GitHub.

\end{itemize}

\section{Objetivos personales}

\begin{itemize}
    \item Crear una herramienta que sea de utilidad para la comunidad universitaria.

    \item Aplicar conocimientos adquiridos durante el grado, especialmente en el ámbito
          de Desarrollo Avanzado de Sistemas Software y Minería de Datos.

    \item Consolidar conocimientos en las siguientes tecnologías: Python, git, docker.

    \item Mejorar habilidades de diseño de software, aplicación de buenas practicas en
          desarrollo, testing de software y documentación.
\end{itemize}

